%-*-latex-*-
%-*-latex-*-
\newcommand\COURSE{ciss240}
\newcommand\ASSESSMENT{q1601}
\newcommand\ASSESSMENTTYPE{Quiz}
\newcommand\POINTS{\textwhite{xxx/xxx}}

\input{myquizpreamble}
\input{yliow}
\input{\COURSE}
%-*-latex-*-

%-*-latex-*-

%-*-latex-*-
\renewcommand\TITLE{\ASSESSMENTTYPE \ \ASSESSMENT}

\renewcommand\EMAIL{}
\renewcommand\AUTHOR{eagraham1@cougars.ccis.edu}

\textwidth=6in
\begin{document}
\topmatterthree

This is a closed-book, no compiler, 5 minute quiz.

%------------------------------------------------------------------------------
\nextq
What is the output of this code fragment:
\begin{console}[fontsize=\scriptsize]
int i = 5, j = 3;
std::cout << "start\n";
if (i > 2)
{
    std::cout << i << ' ' << j << std::endl;
    int j = 6;
    std::cout << i << ' ' << j << std::endl;
    if (j < 5)
    {
        int i = 8;
        std::cout << i << ' ' << j << std::endl;
    }
    else
    {
        int i = 9;
        std::cout << i << ' ' << j << std::endl;
    }
    std::cout << i << ' ' << j << std::endl;
}
else
{
    std::cout << i << ' ' << j << std::endl;
    int j = 7;
    std::cout << i << ' ' << j << std::endl;
    if (j < 5)
    {
        int i = 6;
        std::cout << i << ' ' << j << std::endl;
    }
    else
    {
        int i = 7;
        std::cout << i << ' ' << j << std::endl;
    }
    std::cout << i << ' ' << j << std::endl;
}
std::cout << i << ' ' << j << std::endl;
std::cout << "end\n";
\end{console}
\ANSWER
\begin{answercode}
start
5 3
5 6
9 6
5 6
5 3
end
\end{answercode}

%------------------------------------------------------------------------------
\nextq
In the following code fragment
\begin{console}[fontsize=\scriptsize, numbers=left]
int i = 5;
int i = 6;
\end{console}
there's an error at line 2 because there's a re-declaration of \verb!i!.

For the code fragment below,
write down the line numbers of every statements that has an error.
They all involve variable declaration.
If there's an error for lines 4, 5, 6, 8 you write \verb!4 5 6 8!
with the line numbers in ascending order and a space between two consecutive
line numbers.
\begin{console}[fontsize=\scriptsize, numbers=left]
int i = 5, j = 3;

if (i > 2)
{
    std::cout << i << ' ' << j << std::endl;
    int j = 6;
    std::cout << i << ' ' << j << std::endl;
    if (j < 5)
    {
         int i = 8;
         int j = 7;
         std::cout << i << ' ' << j << std::endl;
    }
    else
    {
        int i = 9;
        int j = 10;
        std::cout << i << ' ' << j << std::endl;
    }
    int j = 11;
    std::cout << i << ' ' << j << std::endl;
}
else
{
    std::cout << i << ' ' << j << std::endl;
    int j = 7;
    std::cout << i << ' ' << j << std::endl;
    if (j < 5)
    {
        int i = 6;
        int j = 7;
        std::cout << i << ' ' << j << std::endl;
    }
    else
    {
        int i = 7;
        int j = 8;
        std::cout << i << ' ' << j << std::endl;
    }
    int j = 9;
    std::cout << i << ' ' << j << std::endl;
}
int j = 10;
std::cout << i << ' ' << j << std::endl;
\end{console}
\ANSWER
\begin{answercode}
5 10 24 30 44 47
\end{answercode}

%------------------------------------------------------------------------------
%-*-latex-*-
\newpage
\input{instructions}
\end{document}

